% BHU PYQ -- Transcribed & Corrected
% Paper: BHUVA12: Environmental Science (Mid Term Test)
% Exam: NEP Value Added Course Environmental Science Mid Term Test, Semester I Examination 2024-25
% Department: Faculty of Science, Banaras Hindu University

\documentclass[11pt,a4paper]{article}
\usepackage[a4paper,top=2.5cm,bottom=2.5cm,left=2.8cm,right=2.8cm,headheight=14pt]{geometry}
\usepackage{amsmath,amssymb}
\usepackage[T1]{fontenc}
\usepackage[utf8]{inputenc}
\usepackage{lmodern,microtype,fancyhdr,enumitem,hyperref,lastpage,parskip}

\hypersetup{
    colorlinks=true,
    urlcolor=black,
    linkcolor=black,
    pdftitle={BHUVA12: Environmental Science Mid Term Test -- BHU Sem I 2024-25 (NEP VAC)}
}

\pagestyle{fancy}
\fancyhf{}
\renewcommand{\headrulewidth}{0.4pt}
\renewcommand{\footrulewidth}{0.4pt}
\fancyhead[L]{\small \textit{Banaras Hindu University}}
\fancyhead[R]{\small \textit{BHUVA12: Environmental Science (Sem I Mid Term 2024-25)}}
\fancyfoot[C]{\small Page \thepage\ of \pageref{LastPage}}

\newlist{parts}{enumerate}{2}
\setlist[parts,1]{label=(\alph*),leftmargin=*,itemsep=2pt,topsep=2pt}

\newcommand{\pts}[1]{\hfill \textit{[#1]}}

\begin{document}

\begin{center}
  {\large\bfseries BANARAS HINDU UNIVERSITY}\\[2pt]
  {\normalsize Faculty of Science}\\[2pt]
  {\normalsize Semester I Examination 2024-25 (NEP)}\\[6pt]
  \rule{0.8\linewidth}{0.5pt}\\[6pt]
  {\large\bfseries Value Added Course (VAC): Environmental Science}\\[4pt]
  {\large\bfseries Mid Term Test --- Paper Code: BHUVA12}\\[4pt]
  {\normalsize Time: One Hour \hfill Maximum Marks: 30}
\end{center}

\rule{\linewidth}{0.4pt}

\smallskip

\noindent \textit{Instructions: The question-booklet contains 30 objective type questions. Candidates must attempt all questions. Each question has four options. Darken any one correct answer. If more than one circle is darked answer will not be evaluated and zero marks will be awarded. Use of log tables or calculators is not permitted.}

\bigskip

\noindent \textbf{Question 1.} Which component of an ecosystem includes all the physical conditions that influence the life of organism within it?\pts{1}
\begin{parts}
  \item Abiotic component
  \item Biotic component
  \item Chemical component
  \item Genetic component
\end{parts}

\medskip

\noindent \textbf{Question 2.} Who proposed the term 'Ecosystem'?\pts{1}
\begin{parts}
  \item A.G. Tansley (1935)
  \item H. Reiter (1868)
  \item Ernst Haeckel (1886)
  \item Grinnel (1971)
\end{parts}

\medskip

\noindent \textbf{Question 3.} In a food chain, herbivores are --\pts{1}
\begin{parts}
  \item Primary producer
  \item Secondary consumer
  \item Primary consumer
  \item Tertiary consumer
\end{parts}

\medskip

\noindent \textbf{Question 4.} River is an example of --\pts{1}
\begin{parts}
  \item Lentic ecosystem
  \item Lotic ecosystem
  \item Marine ecosystem
  \item Terrestrial ecosystem
\end{parts}

\medskip

\noindent \textbf{Question 5.} Phytoplanktons are the component of the which following ecosystem --\pts{1}
\begin{parts}
  \item Forest
  \item Grassland
  \item Desert
  \item Pond
\end{parts}

\medskip

\noindent \textbf{Question 6.} Which of the following is not an ecosystem service?\pts{1}
\begin{parts}
  \item Soil quality maintenance
  \item Decomposition of waste
  \item Weathering of rock
  \item Climate regulation
\end{parts}

\medskip

\noindent \textbf{Question 7.} Which of the following is not a producer?\pts{1}
\begin{parts}
  \item Plant
  \item Herbivore
  \item Chemosynthetic microorganism
  \item Grass
\end{parts}

\medskip

\noindent \textbf{Question 8.} Which of the following statement is not true?\pts{1}
\begin{parts}
  \item Flow of energy in ecosystem is always unidirectional
  \item Detritus food chain starts with dead organic matter
  \item Chemosynthetic microorganisms are producer
  \item Secondary productivity is related to decomposer
\end{parts}

\medskip

\noindent \textbf{Question 9.} Which one of the following is considered as the major cause of biodiversity loss?\pts{1}
\begin{parts}
  \item Climate Change
  \item Global Warming
  \item Habitat Loss
  \item Invasive species
\end{parts}

\medskip

\noindent \textbf{Question 10.} Which one of the following is not an example of \textit{in-situ} conservation strategies?\pts{1}
\begin{parts}
  \item National Park
  \item Botanical Garden
  \item Wildlife Sanctuary
  \item Biosphere Reserve
\end{parts}

\medskip

\noindent \textbf{Question 11.} Which one of the following is an example of national park?\pts{1}
\begin{parts}
  \item Kaziranga
  \item Palamau
  \item Pakke
  \item Sessa
\end{parts}

\medskip

\noindent \textbf{Question 12.} Highest biodiversity used to be found in --\pts{1}
\begin{parts}
  \item Tropical semi-evergreen forests
  \item Tropical rain forests
  \item Tropical deciduous forests
  \item Tropical Grassland
\end{parts}

\medskip

\noindent \textbf{Question 13.} Which of the following is not a level of biodiversity?\pts{1}
\begin{parts}
  \item Genetic
  \item Ecosystem
  \item Gamma
  \item Species
\end{parts}

\medskip

\noindent \textbf{Question 14.} Which of the following statement is true?\pts{1}
\begin{parts}
  \item \textit{In-Situ} conservation is about conservation of biodiversity in their natural habitats
  \item Human overpopulation is the prime cause of biodiversity loss
  \item Overexploitation is about the exploitation of biodiversity for fulfilling needs
  \item Botanical garden is an example of in-situ conservation strategy
\end{parts}

\medskip

\noindent \textbf{Question 15.} Which of the following is not true for invasive species?\pts{1}
\begin{parts}
  \item Invasive species are non-native in origin
  \item Invasive species out-compete our native species
  \item Invasive species are the second-largest cause of biodiversity loss
  \item Invasive species are ecologically important
\end{parts}

\medskip

\noindent \textbf{Question 16.} In the earth's crust, which component stores the most freshwater?\pts{1}
\begin{parts}
  \item Lakes/River/Pond
  \item Groundwater
  \item Oceans
  \item Glaciers
\end{parts}

\medskip

\noindent \textbf{Question 17.} What is the incidence of landmass movement and rockfalls?\pts{1}
\begin{parts}
  \item Landslide
  \item Earth progression
  \item Scaling off
  \item Weathering
\end{parts}

\medskip

\noindent \textbf{Question 18.} San Andreas Fault is an example of --\pts{1}
\begin{parts}
  \item Convergent Plate Boundary
  \item Transform Plate Boundary
  \item Divergent Plate Boundary
  \item American Plate Boundary
\end{parts}

\medskip

\noindent \textbf{Question 19.} In natural conditions, which process is responsible for the breakdown of rocks \textit{in situ}?\pts{1}
\begin{parts}
  \item Erosion
  \item Weathering
  \item Mass wasting
  \item Degradation
\end{parts}

\medskip

\noindent \textbf{Question 20.} What is the name of the type of rock that forms under high pressure and temperature conditions?\pts{1}
\begin{parts}
  \item Igneous
  \item Sedimentary
  \item Metamorphic
  \item Basalt and Granite
\end{parts}

\medskip

\noindent \textbf{Question 21.} Which layer of the atmosphere is closest to Earth's surface?\pts{1}
\begin{parts}
  \item Stratosphere
  \item Mesosphere
  \item Exosphere
  \item Thermosphere
\end{parts}

\medskip

\noindent \textbf{Question 22.} There is a layer of the Earth's surface that is referred to as the \underline{\hspace{2cm}}, which is a part of the lithosphere?\pts{1}
\begin{parts}
  \item Inner Core
  \item Upper Mantle and Crust
  \item Troposphere
  \item Hydrosphere
\end{parts}

\medskip

\noindent \textbf{Question 23.} Which gas of the following is the most abundant in the Earth's atmosphere?\pts{1}
\begin{parts}
  \item Helium
  \item Hydrogen
  \item Nitrogen
  \item Argon
\end{parts}

\medskip

\noindent \textbf{Question 24.} During the summer months, the ozone layer acts as a shield against the \underline{\hspace{2cm}} rays which are directed at our planet.\pts{1}
\begin{parts}
  \item Infrared rays
  \item Gamma rays
  \item Ultraviolet rays
  \item More than one of the above
\end{parts}

\medskip

\noindent \textbf{Question 25.} The uppermost layer of the Earth's atmosphere is \underline{\hspace{2cm}}.\pts{1}
\begin{parts}
  \item Thermosphere
  \item Mesosphere
  \item Stratosphere
  \item Exosphere
\end{parts}

\medskip

\noindent \textbf{Question 26.} Hydrosphere includes \underline{\hspace{2cm}}.\pts{1}
\begin{parts}
  \item Water on the surface of the plant
  \item Underground water
  \item Water in the air
  \item All of the above
\end{parts}

\medskip

\noindent \textbf{Question 27.} The intensity of earthquakes is measured on \underline{\hspace{2cm}}.\pts{1}
\begin{parts}
  \item Beaufort scale
  \item Richter scale
  \item Secant scale
  \item Mercalli scale
\end{parts}

\medskip

\noindent \textbf{Question 28.} Sandstone is a \underline{\hspace{2cm}} rock.\pts{1}
\begin{parts}
  \item Igneous Rock
  \item Metamorphic Rock
  \item Sedimentary Rock
  \item All of the above
\end{parts}

\medskip

\noindent \textbf{Question 29.} Which metal is most abundant in the earth's crust?\pts{1}
\begin{parts}
  \item Aluminium
  \item Magnesium
  \item Calcium
  \item Iron
\end{parts}

\medskip

\noindent \textbf{Question 30.} Molten material below the surface of the earth's crust is called \underline{\hspace{2cm}}.\pts{1}
\begin{parts}
  \item Basalt
  \item Magma
  \item Lava
  \item Limestone
\end{parts}

\vfill
\rule{\linewidth}{0.4pt}

\begin{center}\small
\href{https://bhu-exam.vercel.app/}{Click here to visit our website}\\[4pt]
\href{https://me-aryan.vercel.app/}{Click here to visit our contributor}
\end{center}

\end{document}
